% Editorial bootstrap from the immutable source revision.
% This file is canonical and may be edited independently.

\part{Scope, Status, and the Common Fabius--Rvachev Platform}

\section{The shared analytic and probabilistic platform}
\label{sec:common-platform}
\label{is:p1:sec:platform}

The five source reports use the same bounded Fabius function, Rvachev density, random series, Fourier product, and dyadic spline approximants.  They are fixed once in this part; later parts introduce only the notation peculiar to their question.


% Source fragment: Inverse and sampling frontiers; prefix is.
\subsection{Fabius, Rvachev, and the random series}

Let $V_1,V_2,\ldots$ be independent uniform random variables on $[0,1]$, put
$U_k:=2V_k-1$, and define the uncentered and centered laws
\begin{equation}
 X:=\sum_{k=1}^{\infty}2^{-k}V_k,
 \qquad
 Y:=\sum_{k=1}^{\infty}2^{-k}U_k,
 \qquad
 S_n:=\sum_{k=1}^{n}2^{-k}U_k.
 \label{is:p1:eq:Y-prefix}
\end{equation}
Thus $Y=2X-1$.  Throughout the volume, $X$ and $V_k$ refer to the
uncentered $[0,1]$ model, while $Y$ and $U_k$ refer to its centered
$[-1,1]$ model.
The density of $Y$ is Rvachev's even up-function $\RvachevUp$, supported on $[-1,1]$ \cite{is:p1:Rvachev1990,is:p1:Arias1982,is:p1:ProveItPrimary}.
The bounded Fabius function, originating in Fabius's probabilistic construction \cite{is:p1:Fabius1966}, satisfies
\begin{equation}
 F(x)=\RvachevUp(x-1),\qquad 0\le x\le1,
 \label{is:p1:eq:F-up-shift}
\end{equation}
and is a strictly increasing $C^\infty$ bijection from $[0,1]$ to itself.  We write
\begin{equation}
 \FabiusQuantile:=F^{-1}:[0,1]\longrightarrow[0,1].
 \label{is:p1:eq:G-def}
\end{equation}
The reflection laws are
\begin{equation}
 F(1-x)=1-F(x),\qquad \FabiusQuantile(1-y)=1-\FabiusQuantile(y).
 \label{is:p1:eq:reflection}
\end{equation}

Let $p_n$ be the density of $S_n$ and set
\begin{equation}
 \Fn(x):=p_n(x-1),\qquad 0\le x\le1.
 \label{is:p1:eq:Pn-def}
\end{equation}
On every compact subinterval of $(0,1)$, $\Fn$ is strictly increasing for all sufficiently
large $n$.  Its increasing inverse branch will be denoted
\begin{equation}
 \Gn:=(\Fn)^{-1}.
 \label{is:p1:eq:Gn-def}
\end{equation}
The omitted random tail is a scaled copy of $Y$:
\begin{equation}
 Y-S_n\stackrel{d}=2^{-n}Y.
 \label{is:p1:eq:tail-copy}
\end{equation}

\subsection{The finite Thue--Morse spline}

Put $\ThueMorseSignSymbol_j=(-1)^{\BinaryDigitSum(j)}$, the signed Thue--Morse sequence \cite{is:p1:AlloucheShallit}.  The unequal-box inclusion--exclusion formula for the corresponding box spline \cite{is:p1:deBoorSplines} gives
\begin{equation}
 \boxed{
 \Fn(x)=\frac{2^{\binom n2}}{(n-1)!}
 \sum_{j=0}^{2^n-1}\ThueMorseSignSymbol_j
 \PositivePart{x-2^{-n}-j2^{1-n}}^{n-1}.}
 \label{is:p1:eq:finite-TM-Pn}
\end{equation}
Thus every $\Gn(y)$ is the inverse of an explicitly known piecewise polynomial.  Formula
\eqref{is:p1:eq:finite-TM-Pn} alone does not reveal which polynomial piece controls a prescribed
quantile or why its coefficients should simplify.  The local theorem in \cref{is:p1:sec:exact-local}
answers both questions at the inverse-dyadic outputs $y=F(2^{-r})$.

\subsection{Fourier product and reciprocal-sinc coefficients}

With the angular-frequency convention
$\FourierAngular{f}(\xi)=\int f(x)\EulerE^{-\ImaginaryUnit x\xi}\Differential x$,
\begin{equation}
 \begin{aligned}
 \PhiAng(\xi):=\FourierAngular{\RvachevUp}(\xi)
   &=\prod_{k=1}^{\infty}\SincRad(2^{-k}\xi),\\
 \FourierAngular{p_n}(\xi)
   &=\prod_{k=1}^{n}\SincRad(2^{-k}\xi),
 &\SincRad t&:=\frac{\sin t}{t},\qquad \SincRad0:=1.
 \end{aligned}
 \label{is:p1:eq:Phi-products}
\end{equation}
Near the origin define rational coefficients $b_j$ by
\begin{equation}
 \frac1{\PhiAng(\xi)}=\sum_{j=0}^{\infty}b_j\xi^{2j},
 \qquad b_0=1.
 \label{is:p1:eq:b-def}
\end{equation}
Euler's product gives
\begin{equation}
 \log\frac1{\PhiAng(\xi)}=\sum_{k=1}^{\infty}\gamma_k\xi^{2k},
 \qquad
 \gamma_k=\frac{\zeta(2k)}{k\pi^{2k}(4^k-1)}.
 \label{is:p1:eq:gamma-def}
\end{equation}
Equivalently, in Bernoulli form,
\begin{equation}
 \gamma_k=\frac{(-1)^{k+1}2^{2k-1}B_{2k}}
 {k(2k)!(4^k-1)}.
 \label{is:p1:eq:gamma-Bernoulli}
\end{equation}
The coefficient array is generated by complete Bell polynomials \cite{is:p1:Comtet}:
\begin{equation}
 \boxed{
 b_j=\frac1{j!}\,
 \mathcal B_j\bigl(1!\gamma_1,2!\gamma_2,\ldots,j!\gamma_j\bigr),}
 \label{is:p1:eq:b-Bell}
\end{equation}
or by the recurrence
\begin{equation}
 b_j=\frac1j\sum_{k=1}^{j}k\gamma_kb_{j-k}.
 \label{is:p1:eq:b-rec}
\end{equation}
The first values are
\begin{equation}
 b_1=\frac1{18},\qquad
 b_2=\frac{31}{16200},\qquad
 b_3=\frac{2347}{42865200},\qquad
 b_4=\frac{1904369}{1311675120000}.
 \label{is:p1:eq:b-first}
\end{equation}

The merged volume keeps the four polynomial/function families visually
distinct.  The notation used below is summarized here once:
\begin{center}
\begin{tabular}{@{}ll@{}}
\toprule
Object & Reserved notation \\
\midrule
finite-prefix spline and its local inverse
  & $\FiniteSpline{n}$ and $\FiniteSplineInverse{n}$ \\
forward and inverse compositional iterates
  & $G_n=F^{\comp n}$ and $K_n=(F^{-1})^{\comp n}$ \\
uncentered and centered Appell families
  & $P_n$ and $A_n$ \\
angular and $2\pi$ Fourier products
  & $\PhiAng(\xi)$ and $\PhiCyc(z)=\PhiAng(2\pi z)$ \\
endpoint reversion kernel
  & $\EndpointKernel(z,w)$ \\
\bottomrule
\end{tabular}
\end{center}
Thus a superscript ``$\mathrm{spl}$'' always signals a finite convolution
prefix; bare $G_n,K_n$ are reserved for composition iterates; and the two
Fourier coordinates are never silently interchanged.

The corpus proves the uniform derivative expansion
\begin{equation}
 (\Fn)^{(m)}(x)=
 \sum_{j=0}^{J-1}(-1)^jb_j4^{-jn}F^{(m+2j)}(x)
 +\BigO_{m,J}(4^{-Jn})
 \label{is:p1:eq:forward-prefix-expansion}
\end{equation}
uniformly on $[0,1]$ for every fixed $m,J$.  We shall invert this expansion.

\subsection{Terminating jets at inverse powers of two}

Let
\begin{equation}
 x_r:=2^{-r},\qquad y_r:=F(x_r).
 \label{is:p1:eq:xr-yr}
\end{equation}
The global dilation identity implies
\begin{equation}
 F^{(k)}(x_r)=2^{\binom{k+1}{2}}F(2^{k-r}),
 \quad 0\le k\le r,
 \qquad
 F^{(k)}(x_r)=0\quad(k>r).
 \label{is:p1:eq:dyadic-jet}
\end{equation}
Write $d_0=1$ and
\begin{equation}
 d_r=r!\,2^{\binom r2}F(2^{-r}).
 \label{is:p1:eq:d-def}
\end{equation}
Then the nonconstant part of the terminating Taylor jet is
\begin{equation}
 J_r(z):=\sum_{k=1}^{r}\alpha_{r,k}z^k,
 \qquad
 \alpha_{r,k}
 =2^{\binom{k+1}{2}-\binom{r-k}{2}}
 \frac{d_{r-k}}{k!(r-k)!}.
 \label{is:p1:eq:Jr}
\end{equation}
Thus the Taylor series of $F$ at $x_r$ is the polynomial $y_r+J_r(z)$, although
$F(x_r+z)$ is not equal to that polynomial on any neighborhood.  The discrepancy is flat
at $z=0$.

% Source fragment: Inverse iterates; prefix ii.
\section{Fabius, Rvachev, Thue--Morse, and the derivative atlas}
\label{ii:sec:background}

\subsection{The Rvachev bump and the signed global extension}

Let $\RvachevUp$ denote Rvachev's compactly supported $C^\infty$ function, normalized by
\[
 \SupportOperator\RvachevUp=[-1,1],\qquad \RvachevUp(0)=1,
\]
and related to the bounded Fabius function by
\begin{equation}
 \RvachevUp(t)=F(1-\AbsoluteValue t)\qquad(-1\leq t\leq1).
 \label{ii:eq:up-F}
\end{equation}
In particular, $\RvachevUp$ is even, strictly decreasing on $[0,1]$, and
\begin{equation}
 \RvachevUp(\half)=F(\half)=\half.
 \label{ii:eq:up-half}
\end{equation}

For $k\geq0$, put
\begin{equation}
  \eps_k=(-1)^{\BinaryDigitSum(k)},
  \label{ii:eq:tm-sign}
\end{equation}
where $\BinaryDigitSum(k)$ is the sum of the binary digits of $k$.  The signed global extension
$\ThetaF$ can be written blockwise as
\begin{equation}
 \ThetaF(t)=\eps_k\,\RvachevUp\bigl(t-(2k+1)\bigr),
 \qquad 2k\leq t\leq2k+2.
 \label{ii:eq:theta-block}
\end{equation}
It agrees with $F$ on $[0,1]$, satisfies $\ThetaF'(t)=2\ThetaF(2t)$ globally, and obeys
\begin{equation}
 \AbsoluteValue{\ThetaF(t)}\leq1.
 \label{ii:eq:theta-bound}
\end{equation}
The translate signs in \eqref{ii:eq:theta-block} are precisely the Thue--Morse signs.

\subsection{Exact all-order derivatives}

Define the superfactorial dyadic scale
\begin{equation}
  Q_m:=2^{\binom{m+1}{2}}=2^{m(m+1)/2}.
  \label{ii:eq:Qm}
\end{equation}
Repeated differentiation of the dilation equation gives, for every $m\geq1$ and
$x\in(0,1)$,
\begin{equation}
  F^{(m)}(x)=Q_m\ThetaF(2^m x).
  \label{ii:eq:derivative-atlas}
\end{equation}
Consequently
\begin{equation}
  \AbsoluteValue{F^{(m)}(x)}\leq Q_m.
  \label{ii:eq:F-derivative-bound}
\end{equation}
At the matched dyadic grid this becomes the exact Thue--Morse jet law
\begin{equation}
 F^{(m)}\!\left(\frac{a}{2^m}\right)
 =\begin{cases}
 0,&a\text{ even},\\
 Q_m\eps_{\Floor{a/2}},&a\text{ odd}.
 \end{cases}
 \label{ii:eq:dyadic-jet}
\end{equation}
The factor $Q_m$ is the source of nonanalytic growth.  Since
\begin{equation}
  Q_m^{1/m}=2^{(m+1)/2},
  \label{ii:eq:Q-root}
\end{equation}
no fixed ordinary exponential can neutralize it after division by $m!$.
The only serious issue is cancellation among the many chain-rule terms created by
composition.

\subsection{Shape and orbit dynamics}

The repository proves that $F$ is a strictly increasing bijection of $[0,1]$, that
$F'$ is strictly increasing on $[0,\half]$ and strictly decreasing on
$[\half,1]$, and that
\begin{equation}
 F'(0)=0,\qquad F'(\half)=2,\qquad F'(1)=0.
 \label{ii:eq:slope-values}
\end{equation}
Moreover $F$ is strictly convex on $[0,\half]$ and strictly concave on
$[\half,1]$.  Comparing with the chord joining $(0,0)$ and $(\half,\half)$ gives
\begin{equation}
  F(x)<x\quad(0<x<\half),
  \qquad
  F(x)>x\quad(\half<x<1).
  \label{ii:eq:drift}
\end{equation}
Thus every noncentral interior orbit moves monotonically toward an endpoint:
\begin{align}
 0<x<\half&\implies x>F(x)>F^{\comp2}(x)>\cdots>0,
 \label{ii:eq:left-orbit}\\
 \half<x<1&\implies x<F(x)<F^{\comp2}(x)<\cdots<1.
 \label{ii:eq:right-orbit}
\end{align}
This elementary dynamical fact will force strict unimodality of the exponential
weights in the derivative spine sum.
